\section{Experiments}
\label{sec:experiments}

We ask three questions. \textbf{Q1}: can the language express real robots faithfully
(\cref{sec:fidelity})? \textbf{Q2}: do coding agents build valid robots more often with it than
with existing interfaces, and which parts of it matter (\cref{sec:agents})? \textbf{Q3}: are the
models usable for learning without editing (\cref{sec:rl})?

\subsection{Q1: Rebuilding real robots}
\label{sec:fidelity}
Each of the six reference designs was built by matching a public reference model through rounds of
implement, critique, improve and rebuild; when a robot needed something the API could not express
(closed loops for the gripper, thrust for the quadrotor, torque motors for the quadruped), the
library was extended rather than the design worked around. \Cref{tab:designs} is generated from
the build reports.

\begin{table*}[t]
  \centering
  \caption{The bundled designs, from their build reports. LoC: non-blank, non-comment lines of the
  design file. $\Delta m$: total mass against the reference model (omitted where there is no
  independent reference; the G1's link masses are taken from its reference as budgets). Coll.:
  overlapping part pairs at home. Full range: joints whose sweep reaches both declared limits;
  the limits found are physical self-contacts (gripper pads meeting at closure, left and right
  feet meeting at full abduction, the SO-101 elbow stopping at 2.36\,rad as the reference's
  does). Test: the class test of \cref{sec:report}; Microduck's reference model also fails the
  actuator-only stand test at its real mass.}
  \label{tab:designs}
  % Generated by paper/scripts/collect_results.py. Do not edit by hand.
\begin{tabular}{llrrrcrl}
\toprule
Design & Type & LoC & Mass & $\Delta m$ & Coll. & Full range & Test \\
\midrule
\texttt{so101} & arm & 112 & 599.6\,g & -5.1\% & 0 & 5/6 & arm: pass \\
\texttt{robotiq\_2f85} & gripper & 62 & 1055.6\,g & +0.2\% & 0 & 0/1 & -- \\
\texttt{crazyflie2} & quadrotor & 63 & 27.1\,g & +0.4\% & 0 & -- & hover: pass \\
\texttt{microduck} & biped & 242 & 691.6\,g & -6.2\% & 0 & 8/14 & stand: fail \\
\texttt{microduck\_jetpack} & biped + fans & 141 & 1247.6\,g & -- & 0 & 8/16 & hover: pass \\
\texttt{go2} & quadruped & 125 & 14.85\,kg & -2.4\% & 0 & 6/12 & stand: pass \\
\texttt{g1} & humanoid & 69 & 33.34\,kg & -- & 0 & 15/29 & stand: pass \\
\texttt{duck\_leg} & leg & 58 & 139.3\,g & -- & 0 & 5/5 & -- \\
\bottomrule
\end{tabular}

\end{table*}

\planned{Joint-axis agreement per design from \code{scripts/compare\_joints.py} against each
reference (max distance between axis lines in mm, max angle in degrees) as two extra columns;
README reports $\le$0.2\,mm for SO-101, $\le$0.8\,mm for Microduck, coincident for the others.
Extend \code{collect\_results.py} to run it.}
\planned{Authoring effort: lines of the reference MJCF/URDF vs.\ lines of the design file.}

\subsection{Q2: Coding agents designing robots}
\label{sec:agents}
\planned{Not run yet. Design below; see \code{paper/README.md} for the protocol.}

\paragraph{Tasks.} (i) \emph{Rebuild}: given a reference's published numbers (joint table,
masses, dimensions) but not its model file, produce a model. (ii) \emph{Specify}: given a
natural-language brief (e.g.\ a 4-DoF arm reaching 300\,mm holding 100\,g), produce a design that
passes the class test. (iii) \emph{Repair}: given a design with injected faults (shifted
connector, flipped axis, oversized parent blocking a joint), fix it.

\paragraph{Conditions.} The same agent and budget with: (A) raw MJCF/URDF authoring with MuJoCo
available; (B) build123d with free placement plus a script that exports MJCF; (C) \ms; and
ablations of (C) without the contact check, without the sweep, without the class test, and
without declared parameter ranges.

\paragraph{Metrics.} Rate of valid models (loads, connected, no home collisions), joint-axis error
to the reference, mass error, class-test pass rate, iterations and tokens to first valid model.
An independent checker computes the metrics for every condition, so (A) and (B) are scored by the
same contact and sweep code that \ms uses.

\subsection{Q3: Training on generated models}
\label{sec:rl}
The MJX export of \code{microduck} trains with Brax PPO~\citep{freeman2021brax,schulman2017ppo}
to track 0.3\,m/s velocity commands with an even gait after $3\times10^8$ environment steps
(about 25 minutes on one GPU), and a single policy for \code{microduck\_jetpack} walks, hops,
flies and lands. Both use the build outputs unedited.
\planned{Report learning curves and the hardware used; state seeds and variance.}
